\documentclass[../portafolio.tex]{subfiles}

% Solo agregue paquetes en el preámbulo de ../portafolio.tex

\begin{document}



%%%%%%%%%%%%%%%%%%%%%%%%%%%%%%%%%%%%%%%%%%%%%%%%%%%%%%%%%%%%%%%%%%%%%%%%%%%%%%%%
\section{Git Version Control}  

%\hfill \textbf{Fecha de la actividad:} 19 de agosto de 2022

\medskip

%---------------------------------------------------------------------------------
% Introducción/objetivos de la actividad
%---------------------------------------------------------------------------------
% Activity Overview

This activity introduced the use of Git and GitHub as tools for version control, project management, and collaborative software development. The objective was to learn the fundamental workflows required to track changes, maintain project history, and manage source code repositories effectively.

%---------------------------------------------------------------------------------
% Work Modality

This activity was completed individually during laboratory sessions.

%---------------------------------------------------------------------------------
% Evidence

\subsection{Version Control and Repository Management with Git}

The laboratory began with an introduction to GitHub and its role in modern software development and scientific computing. A personal repository was created to store and organize the projects developed throughout the course, providing a centralized platform for version control and documentation.

Next, several core Git commands were explored to interact with remote repositories. The command \verb|git clone <repository_url>| was used to create a local copy of a remote repository. This workflow enables developers and researchers to work on projects locally while maintaining synchronization with a central repository.

The commands \verb|git pull| and \verb|git push| were then introduced. These commands allow users to retrieve updates from a remote repository and upload local changes, respectively.

A standard Git workflow was also practiced. First, files were staged using \verb|git add|, which prepares modifications for inclusion in the next commit. Changes were then recorded using \verb|git commit -m "message"|, where descriptive commit messages document the purpose of each modification. Finally, the updated project state was synchronized with the remote repository using \verb|git push|.

Through this workflow, a complete version history of the project can be maintained, facilitating reproducibility, collaboration, and project organization.

\subsection{Conclusion}

This activity provided practical experience with Git and GitHub, two of the most widely used tools in software development, data science, and computational research.

Beyond their usefulness for managing coursework, version control systems play an essential role in professional and academic environments by enabling reproducible workflows, collaborative development, and systematic project documentation. The skills acquired in this laboratory establish a foundation for future work involving scientific programming, computational modeling, and collaborative research projects.




\end{document}
